% preamble-exam-zh.tex — 共享 preamble
% 用于所有 exam-zh 模板
%
% 样式策略：
%   屏幕 → 语义彩色（蓝=关键想法，红=易错，绿=路线，灰=提示/教师备注）
%   打印 → 自动降级为黑白（通过 print mode 或 monochrome 选项）

\documentclass{exam-zh}

% ---------- 基础配置 ----------
\examsetup{
  page/size = a4paper,
  page/show-head = false,
  page/show-foot = false,
  sealline/show = false,
  font = times,
  math-font = xits,
}

\usepackage{tcolorbox}
\usepackage{needspace}
\usepackage{graphicx}
\tcbuselibrary{skins,breakable}

% ---------- 打印/屏幕颜色切换 ----------
% 用法：\documentclass[monochrome]{exam-zh} 即可切换为纯黑白
% 注意：不要使用 \if@xxx 放在 makeatother 外部；这里改成普通条件名。
\newif\ifeduprintcolor
\eduprintcolortrue

\makeatletter
\@ifclasswith{exam-zh}{monochrome}{\eduprintcolorfalse}{}
\makeatother

% ---------- 语义颜色定义 ----------
% 屏幕：有色彩区分；打印：降级为灰度
\ifeduprintcolor
  % --- 屏幕彩色模式 ---
  \definecolor{edu-blue}{RGB}{47, 82, 122}       % 关键想法
  \definecolor{edu-blue-bg}{RGB}{243, 246, 252}
  \definecolor{edu-red}{RGB}{180, 40, 40}         % 易错提醒
  \definecolor{edu-red-bg}{RGB}{255, 245, 240}
  \definecolor{edu-green}{RGB}{34, 100, 60}       % 路线图
  \definecolor{edu-green-bg}{RGB}{242, 250, 245}
  \definecolor{edu-orange}{RGB}{160, 90, 0}       % 训练目标
  \definecolor{edu-orange-bg}{RGB}{255, 248, 235}
  \definecolor{edu-gray}{RGB}{100, 100, 100}      % 提示/教师备注
  \definecolor{edu-gray-bg}{RGB}{248, 248, 248}
  \definecolor{edu-purple}{RGB}{90, 50, 120}      % 思考题
  \definecolor{edu-purple-bg}{RGB}{248, 244, 252}
\else
  % --- 打印黑白模式 ---
  \definecolor{edu-blue}{gray}{0.15}
  \definecolor{edu-blue-bg}{gray}{0.96}
  \definecolor{edu-red}{gray}{0.25}
  \definecolor{edu-red-bg}{gray}{0.97}
  \definecolor{edu-green}{gray}{0.20}
  \definecolor{edu-green-bg}{gray}{0.97}
  \definecolor{edu-orange}{gray}{0.25}
  \definecolor{edu-orange-bg}{gray}{0.98}
  \definecolor{edu-gray}{gray}{0.40}
  \definecolor{edu-gray-bg}{gray}{0.97}
  \definecolor{edu-purple}{gray}{0.30}
  \definecolor{edu-purple-bg}{gray}{0.98}
\fi
\colorlet{edu-diagram-muted}{edu-gray}

% ---------- 自定义教学环境 ----------
% 共性：enhanced + breakable（跨页不断裂）

% 原题卡片 — 强边框，突出显示
\newtcolorbox{problemcard}{
  enhanced, breakable,
  colback=edu-blue-bg,
  colframe=edu-blue,
  boxrule=1.2pt,
  arc=0pt,
  left=3mm, right=3mm, top=3mm, bottom=3mm,
}

% 解题路线图 — 左边线 + 浅底
\newtcolorbox{routebox}{
  enhanced, breakable,
  colback=edu-green-bg,
  colframe=edu-green!30,
  boxrule=0pt,
  borderline west={2.5pt}{0pt}{edu-green},
  left=3mm, right=3mm, top=3mm, bottom=3mm,
}

% 关键想法 — 蓝色左边线
\newtcolorbox{keyideabox}{
  enhanced, breakable,
  colback=edu-blue-bg,
  colframe=edu-blue!30,
  boxrule=0pt,
  borderline west={2.5pt}{0pt}{edu-blue},
  left=3mm, right=3mm, top=3mm, bottom=3mm,
}

% 易错提醒 — 红色左边线
\newtcolorbox{mistakebox}{
  enhanced, breakable,
  colback=edu-red-bg,
  colframe=edu-red!20,
  boxrule=0pt,
  borderline west={3pt}{0pt}{edu-red},
  left=3mm, right=3mm, top=2.5mm, bottom=2.5mm,
}

% 提示 — 灰色虚线左边线
\newtcolorbox{hintbox}{
  enhanced, breakable,
  colback=edu-gray-bg,
  colframe=edu-gray!20,
  boxrule=0pt,
  borderline west={2pt}{0pt}{edu-gray, dashed},
  left=3mm, right=3mm, top=2mm, bottom=2mm,
}

% 思考题 — 紫色虚线左边线
\newtcolorbox{thinkbox}{
  enhanced, breakable,
  colback=edu-purple-bg,
  colframe=edu-purple!20,
  boxrule=0pt,
  borderline west={2pt}{0pt}{edu-purple, dashed},
  left=2.5mm, right=2.5mm, top=2.5mm, bottom=2.5mm,
}

% 教师备注 — 灰色左边线，小字
\newtcolorbox{teachernote}{
  enhanced, breakable,
  colback=edu-gray-bg,
  colframe=edu-gray!15,
  boxrule=0pt,
  borderline west={2pt}{0pt}{edu-gray!60},
  left=2.5mm, right=2.5mm, top=2.5mm, bottom=2.5mm,
  fontupper=\small,
  colupper=black!60,
}

% 训练目标 — 橙色左边线，小字
\newtcolorbox{traininggoal}{
  enhanced, breakable,
  colback=edu-orange-bg,
  colframe=edu-orange!15,
  boxrule=0pt,
  borderline west={1.5pt}{0pt}{edu-orange},
  left=2mm, right=2mm, top=1mm, bottom=1mm,
  fontupper=\small,
}

% 答题区 — 无边框，纯留白
\newcommand{\answerarea}[1][40mm]{%
  \vspace{2mm}%
  \begin{tcolorbox}[
    enhanced, breakable,
    colback=white,
    colframe=white,
    boxrule=0pt,
    height=#1,
    valign=top,
    bottom=0mm,
    after skip=0mm,
  ]
  \end{tcolorbox}%
}

\newcommand{\diagramcolfigure}[3]{%
  \begin{minipage}[t]{#2}
    \vspace{0pt}\centering
    \includegraphics[width=\linewidth]{\detokenize{#1}}
    \if\relax\detokenize{#3}\relax\else
      \par\vspace{1mm}{\scriptsize\color{edu-diagram-muted}#3}
    \fi
  \end{minipage}%
}

\newcommand{\diagramrowitem}[4]{%
  \begin{minipage}[t]{#2}
    \vspace{0pt}\centering
    \includegraphics[width=\linewidth]{\detokenize{#1}}
    \if\relax\detokenize{#3}\relax\else
      \par\vspace{1mm}{\scriptsize\bfseries #3}
    \fi
    \if\relax\detokenize{#4}\relax\else
      \par\vspace{0.5mm}{\scriptsize\color{edu-diagram-muted}#4}
    \fi
  \end{minipage}%
}

\newcommand{\answerareawithdiagramcol}[4]{%
  \vspace{2mm}%
  \begin{tcolorbox}[
    enhanced, breakable,
    colback=white,
    colframe=white,
    boxrule=0pt,
    height=#1,
    valign=top,
    bottom=0mm,
    after skip=0mm,
  ]
  \begin{minipage}[t]{\dimexpr\linewidth-#3-6mm\relax}
    \vspace{0pt}
  \end{minipage}\hfill
  \diagramcolfigure{#2}{#3}{#4}
  \end{tcolorbox}%
}

\newcommand{\answerareawithdiagram}[4]{%
  \answerareawithdiagramcol{#1}{#2}{#3}{#4}%
}

% ---------- 字体设置 ----------
% exam-zh (ctexbook) already sets CJK fonts via ctex.
% Only override if specific fonts are needed.
% macOS: Songti SC, STHeiti, STFangsong
% Linux: SimSun, SimHei, FangSong

% ---------- 页面设置 ----------
% exam-zh (ctexbook) already loads geometry.
% Override margins if needed:
% \geometry{top=16mm, bottom=16mm, left=14mm, right=14mm}

\examsetup{
  question/show-answer = true,
  fillin/show-answer = true,
  solution/show-solution = show-stay,
}

% ---------- 教师版解答样式 ----------
\newtcolorbox{solutionblock}{
  enhanced, breakable,
  colback=edu-blue-bg,
  colframe=edu-blue!25,
  boxrule=0pt,
  borderline west={2pt}{0pt}{edu-blue!50},
  arc=0pt,
  left=3mm, right=3mm, top=2mm, bottom=2mm,
  before skip=2mm,
  after skip=3mm,
  fontupper=\small,
}

% 解答步骤编号
\newcounter{solstep}
\newcommand{\solstep}[2]{%
  \stepcounter{solstep}%
  \par\medskip
  \noindent\hangindent=6mm
  {\color{edu-blue}\bfseries\textcircled{\scriptsize\thesolstep}\enspace #1}\enspace #2\par
}

\begin{document}

\section*{一次函数与反比例函数交点面积综合 · 练习（教师版）}

\section*{练习题（每题 10 分，共 40 分）}

\needspace{8\baselineskip}

\begin{problem}[points=10]
  已知一次函数 $y_1=kx+b$ 与反比例函数 $y_2=\dfrac{m}{x}$ 的图像交于点 $A,B$。已知：点 $A$ 的横坐标为 $3$，点 $B$ 的横坐标为 $-1$，一次函数图像与 $y$ 轴交于点 $C(0,-2)$。
\begin{enumerate}[label=(\arabic*)]
  \item 求一次函数解析式、反比例函数解析式，以及点 $A,B$ 的坐标；
  \item 根据图像求不等式 $kx+b>\dfrac{m}{x}$ 的解集；
  \item 求 $\triangle OAB$ 的面积。
\end{enumerate}

  \answerareawithdiagramcol{58mm}{diagrams/p1.pdf}{62mm}{第 1 题图}
  \setcounter{solstep}{0}
  \begin{solutionblock}
    \textbf{答案：}(1) $y_1=x-2$，$y_2=\dfrac3x$，$A(3,1)$，$B(-1,-3)$；(2) $-1<x<0$ 或 $x>3$；(3) $S_{\triangle OAB}=4$。\par\medskip
    \solstep{求函数和交点}{由 $C(0,-2)$ 得 $b=-2$。设 $A(3,3k-2)$，$B(-1,-k-2)$。由 $3(3k-2)=(-1)(-k-2)$，得 $k=1$，所以 $A(3,1)$，$B(-1,-3)$，$m=3$。}
    \solstep{解不等式}{$x-2>\dfrac3x$ 看直线在双曲线上方。断点为 $-1,0,3$，取 $(-1,0)$ 和 $(3,+\infty)$。}
    \solstep{求面积}{直线 $AB$ 与 $y$ 轴交于 $P(0,-2)$，所以 $OP=2$。把 $\triangle OAB$ 按 $OP$ 拆成 $\triangle OPA$ 和 $\triangle OPB$；两者共用底 $OP$，高分别是 $|x_A|$、$|x_B|$。因 $A,B$ 在 $y$ 轴两侧，$S_{\triangle OAB}=S_{\triangle OPA}+S_{\triangle OPB}=\dfrac12OP(|x_A|+|x_B|)=\dfrac12OP|x_A-x_B|=\dfrac12\cdot2\cdot|3-(-1)|=4$。}
  \end{solutionblock}
\end{problem}


\needspace{8\baselineskip}

\begin{problem}[points=10]
  已知一次函数 $y_1=kx+b$ 与反比例函数 $y_2=\dfrac{m}{x}$ 的图像交于点 $A,B$。已知：$k=-1$，$b=5$，点 $A$ 的横坐标为 $1$。
\begin{enumerate}[label=(\arabic*)]
  \item 求一次函数解析式、反比例函数解析式，以及点 $A,B$ 的坐标；
  \item 根据图像求不等式 $kx+b>\dfrac{m}{x}$ 的解集；
  \item 求 $\triangle OAB$ 的面积。
\end{enumerate}

  \answerareawithdiagramcol{58mm}{diagrams/p2.pdf}{62mm}{第 2 题图}
  \setcounter{solstep}{0}
  \begin{solutionblock}
    \textbf{答案：}(1) $y_1=-x+5$，$y_2=\dfrac4x$，$A(1,4)$，$B(4,1)$；(2) $x<0$ 或 $1<x<4$；(3) $S_{\triangle OAB}=\dfrac{15}{2}$。\par\medskip
    \solstep{求函数和交点}{由 $k=-1,b=5$，得 $y_1=-x+5$。点 $A$ 的横坐标为 $1$，所以 $A(1,4)$，从而 $m=1\cdot4=4$，$y_2=\dfrac4x$。解 $-x+5=\dfrac4x$，得 $x=1$ 或 $x=4$，所以 $B(4,1)$。}
    \solstep{解不等式}{$-x+5>\dfrac4x$ 的断点为 $0,1,4$；直线在上方的区间为 $(-\infty,0)$ 和 $(1,4)$。}
    \solstep{求面积}{直线 $AB$ 与 $y$ 轴交于 $P(0,5)$，所以 $OP=5$。把图形按 $OP$ 拆看，$\triangle OPB$ 和 $\triangle OPA$ 共用底 $OP$，高分别为 $|x_B|$、$|x_A|$。因 $A,B$ 在 $y$ 轴同侧，所求区域是大三角形减小三角形：$S_{\triangle OAB}=S_{\triangle OPB}-S_{\triangle OPA}=\dfrac12OP(|x_B|-|x_A|)=\dfrac12OP|x_A-x_B|=\dfrac12\cdot5\cdot|1-4|=\dfrac{15}{2}$。}
  \end{solutionblock}
\end{problem}


\needspace{8\baselineskip}

\begin{problem}[points=10]
  已知一次函数 $y_1=kx+b$ 与反比例函数 $y_2=\dfrac{m}{x}$ 的图像交于点 $A,B$。已知：$b=1$，$m=6$，点 $A$ 的横坐标为 $2$。
\begin{enumerate}[label=(\arabic*)]
  \item 求一次函数解析式、反比例函数解析式，以及点 $A,B$ 的坐标；
  \item 根据图像求不等式 $kx+b>\dfrac{m}{x}$ 的解集；
  \item 已知 $P(0,4)$，求 $\triangle PAB$ 的面积。
\end{enumerate}

  \answerareawithdiagramcol{62mm}{diagrams/p3.pdf}{62mm}{第 3 题图}
  \setcounter{solstep}{0}
  \begin{solutionblock}
    \textbf{答案：}(1) $y_1=x+1$，$y_2=\dfrac6x$，$A(2,3)$，$B(-3,-2)$；(2) $-3<x<0$ 或 $x>2$；(3) $S_{\triangle PAB}=\dfrac{15}{2}$。\par\medskip
    \solstep{求函数和交点}{由 $m=6$ 且 $x_A=2$，得 $A(2,3)$。又 $b=1$，代入 $y=kx+1$ 得 $2k+1=3$，所以 $k=1$，$y_1=x+1$。解 $x+1=\dfrac6x$，得 $x=2$ 或 $x=-3$，所以 $B(-3,-2)$。}
    \solstep{解不等式}{$x+1>\dfrac6x$ 的断点为 $-3,0,2$；直线在上方的区间为 $(-3,0)$ 和 $(2,+\infty)$。}
    \solstep{求 $PAB$ 面积}{因为题中已有点 $P(0,4)$，把直线 $AB$ 与 $y$ 轴交点记为 $Q(0,1)$，则 $PQ=3$。把 $\triangle PAB$ 按 $PQ$ 拆成 $\triangle PQA$ 和 $\triangle PQB$；两者共用底 $PQ$，高分别是 $|x_A|$、$|x_B|$。因 $A,B$ 在 $y$ 轴两侧，$S_{\triangle PAB}=S_{\triangle PQA}+S_{\triangle PQB}=\dfrac12PQ(|x_A|+|x_B|)=\dfrac12PQ|x_A-x_B|=\dfrac12\cdot3\cdot|2-(-3)|=\dfrac{15}{2}$。}
  \end{solutionblock}
\end{problem}


\needspace{8\baselineskip}

\begin{problem}[points=10]
  已知一次函数 $y_1=kx+b$ 与反比例函数 $y_2=\dfrac{m}{x}$ 的图像交于点 $A,B$。已知：$k=-1$，$m=-2$，点 $A$ 的横坐标为 $1$。
\begin{enumerate}[label=(\arabic*)]
  \item 求一次函数解析式、反比例函数解析式，以及点 $A,B$ 的坐标；
  \item 根据图像求不等式 $kx+b>\dfrac{m}{x}$ 的解集；
  \item 点 $P$ 在反比例函数第四象限分支上，且在点 $A$ 的右侧。若 $S_{\triangle PAB}=3$，求点 $P$ 的坐标。
\end{enumerate}

  \answerareawithdiagramcol{64mm}{diagrams/p4.pdf}{62mm}{第 4 题图}
  \setcounter{solstep}{0}
  \begin{solutionblock}
    \textbf{答案：}(1) $y_1=-x-1$，$y_2=-\dfrac2x$，$A(1,-2)$，$B(-2,1)$；(2) $x<-2$ 或 $0<x<1$；(3) $P(2,-1)$。\par\medskip
    \solstep{求函数和交点}{由 $m=-2$ 且 $x_A=1$，得 $A(1,-2)$。又 $k=-1$，代入 $y=-x+b$ 得 $-1+b=-2$，所以 $b=-1$，$y_1=-x-1$。解 $-x-1=-\dfrac2x$，得 $x=1$ 或 $x=-2$，所以 $B(-2,1)$。}
    \solstep{解不等式}{$-x-1>-\dfrac2x$。由图像可知，直线在双曲线上方的区间为 $(-\infty,-2)$ 和 $(0,1)$。}
    \solstep{设点并列面积方程}{点 $P$ 在第四象限分支且在 $A$ 右侧，设 $P(t,-\dfrac2t)$，其中 $t>1$。直线 $AB$ 为 $y=-x-1$，过 $P$ 作铅垂线交 $AB$ 于 $H$，则 $H(t,-t-1)$。}
    \solstep{解出 $P$}{此时 $A,B$ 都在铅垂线 $PH$ 的左侧，且 $A$ 在 $B,H$ 之间，所以 $\triangle PHB$ 是大三角形，$\triangle PHA$ 是被减去的小三角形。于是 $S_{\triangle PAB}=S_{\triangle PHB}-S_{\triangle PHA}=\dfrac12PH\big((t-x_B)-(t-x_A)\big)=\dfrac12PH|x_A-x_B|$。又 $PH=|(-t-1)-(-\dfrac2t)|=|t+1-\dfrac2t|$，所以 $\dfrac12\cdot3\cdot|t+1-\dfrac2t|=3$。由 $t>1$ 得 $t+1-\dfrac2t>0$，解得 $t=2$，$P(2,-1)$。}
  \end{solutionblock}
\end{problem}


\clearpage
\section*{参考答案}


\begin{keyideabox}
  \textbf{练习题}
  第 1 题：$y_1=x-2$，$y_2=\dfrac3x$，$A(3,1)$，$B(-1,-3)$；解集为 $-1<x<0$ 或 $x>3$；$S_{\triangle OAB}=4$。

第 2 题：$y_1=-x+5$，$y_2=\dfrac4x$，$A(1,4)$，$B(4,1)$；解集为 $x<0$ 或 $1<x<4$；$S_{\triangle OAB}=\dfrac{15}{2}$。

第 3 题：$y_1=x+1$，$y_2=\dfrac6x$，$A(2,3)$，$B(-3,-2)$；解集为 $-3<x<0$ 或 $x>2$；$S_{\triangle PAB}=\dfrac{15}{2}$。

第 4 题：$y_1=-x-1$，$y_2=-\dfrac2x$，$A(1,-2)$，$B(-2,1)$；解集为 $x<-2$ 或 $0<x<1$；$P(2,-1)$。

\end{keyideabox}



\end{document}
